\documentclass[11pt]{article}
\usepackage[margin=0.9in]{geometry}
\usepackage{amsmath,amssymb,amsthm,booktabs,mathrsfs}
\usepackage[T1]{fontenc}
\usepackage{lmodern,microtype}
\usepackage[colorlinks=true,linkcolor=blue,urlcolor=blue]{hyperref}
\newtheorem{theorem}{Theorem}[section]
\newtheorem{lemma}[theorem]{Lemma}
\newtheorem{proposition}[theorem]{Proposition}
\newtheorem{corollary}[theorem]{Corollary}
\theoremstyle{remark}
\newtheorem{remark}[theorem]{Remark}
\newcommand{\norm}[1]{\left\|#1\right\|}
\title{Toward a four-dimensional Yang--Mills construction\\
\large Notebook 25: a smaller representative of the classical gauge orbit}
\author{Research working document}
\date{17 October 2026}
\begin{document}
\maketitle
\begin{abstract}
The contour-gauge stationary representative need not contract
link angles under blocking. Before applying that rooted current
gauge, however, the partial ordinary average of the established
fine classical minimizer has link-angle bound $C/R$, where
$R=L^s$ is the remaining averaging ratio. We prove the same
factor in its strict rooted support norm and quadratic remainder,
with constants uniform in both averaging depths at fixed base
$L$. The proof applies Banach-valued local analytic estimates
to complete output fibers and does not sum over their increasing
number of coordinates. For an actual one-step gauge-invariant
total correction, this representative supplies an analytic
continuation of its stationary-composition germ wherever the
smaller representative lies in the original domain.
The support-coefficient norm and complete inherited fluctuation
integral still require their own domain estimates.
\end{abstract}

This result addresses the angle-domain obstruction identified in
Notebook 22. It uses the ordinary trajectory bound, the explicit
ordinary minimizer, its localized construction, and the local
gauge identity. At one physical step the displayed amplitude
factor is $C/L$; the constant may depend on that fixed base $L$.
The uniform $C/L^s$ statement for larger $s$ concerns the iterated
constraint, and does not identify a new direct-block integration.


\section{A small representative before the rooted current gauge}
\label{sv-section}

Fix the base odd ratio $L$, dimension four, $0<\alpha<1$, the compact
matrix gauge group and compatible finite periodic geometry. Write
\begin{equation}
 M=L^k,\qquad R=L^s,\qquad N=MR,\qquad
 \eta=N^{-1},\qquad h_t=L^t\eta,\qquad h_k=R^{-1}.
 \label{sv-scales}
\end{equation}
Coordinates in this section have final next-coarse spacing one. The
current lattice after $k$ averages has spacing $1/R$. The datum
$W=e^{iw}$ is a group datum, and $w$ is its link angle. All constants
below are uniform in $k,s$ and volume for the fixed base $L$.

Let $A_{\mathrm c,N}(w)$ be the established corrected fixed-gauge
minimizer, and let $T_N$ be the explicit ordinary-representative
transfer. Thus
\begin{equation}
 A_{\mathrm o,N}(w)=T_N(A_{\mathrm c,N}(w)),\qquad
 U_N(w)=e^{i\eta A_{\mathrm o,N}(w)},\qquad
 M_L^{k+s}(U_N(w))=e^{iw},\qquad
 \|A_{\mathrm o,N}(w)\|_{\mathcal N_\alpha}\le C_{\mathrm o}\|w\|_\infty.
 \label{sv-ordinary-minimizer}
\end{equation}
The last bound and a common complex product domain are uniform in the
total depth $k+s$. Superscripts on $M_L$ mean iteration of the fixed-$L$
ordinary average. Define
\begin{equation}
 V_{\mathrm p,k,s}(w)=M_L^k(U_N(w)),\qquad
 a_{\mathrm p,k,s}(w)=\frac1i\log V_{\mathrm p,k,s}(w).
 \label{sv-pre-representative}
\end{equation}
Exact composition gives
\begin{equation}
 M_L^s(V_{\mathrm p,k,s}(w))=e^{iw}.
 \label{sv-iterated-constraint}
\end{equation}
For $s=1$ this is the pre-root-gauge representative used in the physical
one-step construction. For $s>1$ it is a representative of the iterated
constraint \eqref{sv-iterated-constraint}; a separately defined direct
$R$-block average, pivot chart and centered integral are not identified
with that iteration here.

\subsection{The exact amplitude factor}
\label{sv-amplitude-section}

The proved ordinary-reference trajectory estimate is uniform in every
prefix of the averaging sequence. With $U=e^{i\eta A}$ and
$a=\|A\|_\infty$ sufficiently small, it states
\begin{equation}
 \left\|\frac1i\log M_L^t(U)\right\|_\infty
 \le 3K_Vh_ta,\qquad
 K_V=2+\frac{2K_\Gamma}{L-1},\qquad K_\Gamma=8m_\Gamma,
 \label{sv-prefix-estimate}
\end{equation}
where $m_\Gamma$ is a fixed one-step contour-length bound. In that proof
the only accumulated scale sum is
\begin{equation}
 \sum_{q<t}h_q=\frac{h_t-\eta}{L-1}.
 \label{sv-prefix-sum}
\end{equation}
In particular its constants depend on the base $L$, not on the endpoint
ratio $R$. Apply \eqref{sv-prefix-estimate} with $t=k$ and the ordinary
minimizer in \eqref{sv-ordinary-minimizer}. On a smaller fixed ball,
\begin{equation}
 \boxed{\quad
 \|a_{\mathrm p,k,s}(w)\|_\infty
       \le \frac{C_0}{R}\|w\|_\infty,\qquad
 C_0=3K_VC_{\mathrm o}.
 \quad}
 \label{sv-amplitude}
\end{equation}
This is also a complex estimate on the common logarithm branch. When
$k=0$, it reduces to $a_{\mathrm p,0,s}=\eta A_{\mathrm o,N}$ with
$\eta=1/R$, checking the normalization.

The same trajectory proof gives a uniformly holomorphic physical map
\begin{equation}
 P_{k,N}(A)=\frac1{ih_k}\log M_L^k(e^{i\eta A}),\qquad
 \|P_{k,N}(A)\|_\infty\le 3K_V\|A\|_\infty.
 \label{sv-physical-prefix-map}
\end{equation}
Each output reads a fine-field neighborhood of physical radius at most
$c_Lh_k$: the one-step dependence ranges sum geometrically. No physical
adjoint of this ordinary map is used in this argument.

\subsection{Strict supports without counting the output fiber}
\label{sv-support-section}

Support labels denote final unit cells, carrying the variables $w$.
Let $\tau$ be periodic graph Steiner length, allowing auxiliary vertices.
Fix an admissible exponent $\kappa$ from the classical support theorem,
and any fixed cardinality reserve $a\ge0$. For a current lattice of
spacing $h=1/R$, collect all bond components whose initial sites lie in
one final cell $i$ into a Banach fiber $E_{h,i}$ with its maximum norm.
The fiber dimension is of order $R^d$, but the norm is not a sum over
those coordinates. For a strict support family define
\begin{equation}
 \|u\|_{\kappa,a;r}
 =\sup_i\sum_{\varnothing\ne X}
 e^{\kappa\tau(\{i\}\cup X)+a|X|}
 \sup_{\|w_X\|_\infty<r}\|u_{i,X}(w_X)\|_{E_{h,i}}.
 \label{sv-family-norm}
\end{equation}
The output maximum is inside the support sum. Complete fine-field input
families use the established smooth partition and complete
$\mathcal N_\alpha$ cell spaces.

\begin{theorem}[Supported pre-gauge representative]
\label{sv-support-theorem}
There are $r_a>0$ and $C_a<\infty$, independent of $k,s$ and volume,
such that for $0<r\le r_a$ the representative
\eqref{sv-pre-representative} admits absolutely convergent strict
holomorphic families
\begin{equation}
 a_{\mathrm p,k,s}(w)|_i
   =\sum_{\varnothing\ne X}a_{\mathrm p;i,X}(w_X),\qquad
 a_{\mathrm p;i,X}(0)=0,\qquad
 \|a_{\mathrm p,k,s}\|_{\kappa,a;r}\le \frac{C_ar}{R}.
 \label{sv-supported-amplitude}
\end{equation}
There is a uniformly supported linear operator
$H_{\mathrm p,k,s}=D(P_{k,N}T_N)(0)H_N$ such that
\begin{equation}
 \|a_{\mathrm p,k,s}-R^{-1}H_{\mathrm p,k,s}w\|_{\kappa,a;r}
       \le \frac{C_ar^2}{R}.
 \label{sv-supported-remainder}
\end{equation}
Here $H_N$ is the corrected flat minimizing extension. The selected
linear operator is not identified with a separately rooted source
gauge extension.
\end{theorem}

\begin{proof}
The classical support construction supplies
$A_{\mathrm c,N}(w)$ in complete fine-field cells with norm at most
$C_{a,0}r$, uniformly in $N$ and volume, including its linear part and
quadratic remainder. The ordinary transfer $T_N$ is uniformly
holomorphic $\mathcal N_\alpha\to\mathcal N_\alpha$ on a fixed ball,
vanishes at zero, and has fixed final-physical dependence radius.
Indeed its cumulative endpoint gauge is a local bounded analytic
function of the fine potential, its point extension is a fixed smooth
local interpolation, and its fine gauge action is a local analytic jet
map with the apparent inverse-mesh singularity removed before taking
derivatives.

Consequently $F_{k,N}=P_{k,N}\circ T_N$ is uniformly bounded and
holomorphic on a fixed $\mathcal N_\alpha$ ball into the physical bounded
output field, and $F_{k,N}(0)=0$. Its dependence radius in final units is
fixed. For the entire fiber $E_{h,i}$ the union of the individual output
neighborhoods lies in one fixed collar of $i$, even though the number of
components inside the fiber grows with $R$.

Let $\rho$ be a common analytic input radius, $C_{\rm rec}$ the smooth
reconstruction norm, $D_0$ the fixed number of input cell labels in that
collar, and $R_0$ its graph radius. Banach-valued Cauchy and polarization
give for the symmetric Taylor coefficients, including the factor
$1/n!$, the mixed cell-slot bound
\begin{equation}
 v_n\le Kc^n,\qquad
 c=\frac{eC_{\rm rec}D_0e^{\kappa R_0}}{\rho},\qquad n\ge1.
 \label{sv-vertex-majorant}
\end{equation}
The constant $K$ bounds the output in its maximum norm on the analytic
ball. Cauchy is applied to that whole Banach output; it is not applied
to each scalar component followed by a sum over $R^d$ components.
The coefficient norm places the output maximum before the sum over
input cell labels, precisely as required by
\eqref{sv-family-norm}.

For substituted input supports $X_1,\ldots,X_n$ rooted at
$j_1,\ldots,j_n$, the exact gluing inequality is
\begin{equation}
 \tau\left(\{i\}\cup\bigcup_qX_q\right)
 \le \tau(\{i,j_1,\ldots,j_n\})
                    +\sum_q\tau(\{j_q\}\cup X_q).
 \label{sv-root-gluing}
\end{equation}
Also $e^{a|\cup_qX_q|}\le\prod_qe^{a|X_q|}$. The fixed coefficients of
the local map add no variable support. Thus substitution of an input
family of norm $t$ has output norm at most
\begin{equation}
 \sum_{n\ge1}Kc^nt^n=\frac{Kct}{1-ct}\le 2Kct
                  \quad\hbox{if }ct\le\tfrac12.
 \label{sv-support-majorant}
\end{equation}
Choose $r_a$ so $cC_{a,0}r_a\le1/2$, and shrink it to respect every
previous evaluated-branch domain. These choices are independent of
$R$. Finally $a_{\mathrm p,k,s}=h_kF_{k,N}(A_{\mathrm c,N}(w))$ gives
\eqref{sv-supported-amplitude}. Separating the linear terms in the two
analytic series gives \eqref{sv-supported-remainder}. Absolute
convergence justifies rearrangement and identifies the sum with the
actual representative.
\end{proof}

This proof neither uses the nonuniform ordinary-average physical
adjoint nor infers tree locality from a one-weight response bound. The
spatial exponent is in final unit cells; in current-unit distances it
has length scale $R$. Uniformity in $R=L^s$ keeps the base $L$ fixed.

\subsection{One-step physical identification and analytic continuation}
\label{sv-germ-section}

Set $s=1$. The physical construction uses the actual averaged-contour
current gauge
\begin{equation}
 u(y)=1,\qquad u(x)=P_{V_{\mathrm p}}(y,x)
       \quad(x\hbox{ in the block of }y),\qquad
 V_*(W)=V_{\mathrm p}(W)^u.
 \label{sv-root-gauge}
\end{equation}
It has $P_{V_*}=1$, $M_L(V_*)=W$ and, on the common complex branch,
\begin{equation}
 \|\log u/i\|_\infty
   \le C_L\|a_{\mathrm p}\|_\infty\le C'_L\|w\|_\infty.
 \label{sv-root-gauge-size}
\end{equation}
The classical local gauge-invariance theorem applies to both
configurations and this gauge path after a common fixed shrink. Hence
$S_M(V_*)=S_M(V_{\mathrm p})$. The existing minimum comparisons and
common strict-convexity uniqueness identify \eqref{sv-root-gauge}
with the physical centered branch. Gauge-zero is obtained by the
explicit gauge and that comparison, not from stationarity alone.
The gauge in \eqref{sv-root-gauge-size} need not have size $w/L$;
boundary angles of $V_*$ can be of order $w$ even when
\eqref{sv-amplitude} holds. Thus the earlier pure-gauge obstruction
for the rooted representative is preserved.

Let $E$ be a total holomorphic current-field functional on the angle
product ball $\|a\|_\infty<b$, locally gauge invariant when its two
arguments and the gauge lie in their common invariance domains. Then
\begin{equation}
 E_{\mathrm p}(w)=E(a_{\mathrm p,k,1}(w))
 \label{sv-germ-extension}
\end{equation}
is holomorphic for
\begin{equation}
 \|w\|_\infty<\min\{r_{\rm class},Lb/C_0\},
 \label{sv-germ-domain}
\end{equation}
with smaller margins if required. On a sufficiently small overlap
around zero, both original compositions and the gauge path are
admissible, so
\begin{equation}
 E(a_*(w))=E(a_{\mathrm p,k,1}(w)).
 \label{sv-overlap-equality}
\end{equation}
The real gauge identity and holomorphic uniqueness prove this
equality. Therefore \eqref{sv-germ-extension} is an analytic
continuation of the original stationary-composition germ to
\eqref{sv-germ-domain}. It does not assert that $E$ was already defined
at $a_*(w)$ outside its original ball. No coefficientwise identity for
an arbitrary decomposition $E=\sum_HE_H$ is inferred.

If $E$ is an inherited logarithm, its original domain must already
carry a nonvanishing integral and compatible holomorphic logarithm.
The covariance of a possibly vanishing integral alone is insufficient.
For $s>1$, $E(a_{\mathrm p,k,s})$ remains an analytic composition with
the same amplitude improvement, but identifying it with a particular
macrostep stationary germ requires that macrostep's branch and gauge
comparison. A direct-$R$ chart is not supplied here.

For any fixed $R=L^s>C_0$, the representative itself maps an angle ball
of radius $b$ into the strictly smaller current angle ball of radius
$C_0b/R$, provided $b$ is below the uniform classical radius. This
geometric domain inclusion is proved; invariance is needed to identify
the resulting composition with a stationary germ, not for the inclusion
itself. It is not an activity-norm contraction.

\subsection{The activity norm still has a coordinate cost}
\label{sv-activity-scope}

Suppose the inherited connected-carrier activities have scalar
coordinate polydisc radius $b$, at most $D$ coordinates per actual old
cell, and
\begin{equation}
 \sup_i\sum_{H\ni i}e^{a_0|H|}\|E_H\|_b\le\epsilon.
 \label{sv-old-activity-norm}
\end{equation}
Let $\pi$ project those old cells to final cells, preserving connected
images, with fibers bounded by $F_R$. Put $\beta=a_{\rm out}+\mu$.
The established carrier-transport theorem, applied to
\eqref{sv-supported-amplitude}, gives output norm at most
$F_RS_\delta\epsilon$ whenever
\begin{equation}
 \beta+\delta\le\kappa,\qquad
 \beta+D\log\frac1{1-C_ar/(Rb)}\le a_0.
 \label{sv-activity-budget}
\end{equation}
Here $D$ is the actual old-cell scalar-coordinate count, not the output
Banach fiber dimension in Theorem~\ref{sv-support-theorem}.
Nevertheless if the chosen background/Gaussian cell has
$D\asymp R^d$, the small-radius sufficient condition in
\eqref{sv-activity-budget} still has $r/b=O(R/D)=O(R^{1-d})$.
If $D$ stays fixed, that particular loss disappears but the projection
factor $F_R$ can grow like $R^d$. Analytic-domain inclusion and a
contracting activity-norm induction are distinct assertions.

There is a further pointwise identity on the original common domain:
with $V_*=V_{\mathrm p}^u$,
\begin{equation}
 (e^{iB}V_*)^{u^{-1}}(b)
   =e^{i\operatorname{Ad}_{u(b_-)^{-1}}B(b)}V_{\mathrm p}(b).
 \label{sv-fluctuation-gauge}
\end{equation}
For real $u$ the adjoint preserves the invariant color norm; it is
uniformly bounded on the small complex domain. Using the right side
to extend an inherited insertion with $B=\Phi_{V_*}(gz)$ is a possible
next step. The coupling is inside the nonlinear chart; generally
$\Phi_V(gz)\ne g\Phi_V(z)$. This section does not prove a larger-domain normalized
Gaussian-integral identity. That needs the actual chart and
fluctuation bounds, holomorphic domination, interaction norm control
and equality of the two analytic integrals on an overlap, with the
exact Gaussian, gauge penalty, pivot/Haar factors and cutoff arguments
retained. No stable quantum iteration or running-coupling conclusion
is made.


\begin{thebibliography}{9}
\bibitem{prior25}
\emph{Toward a four-dimensional Yang--Mills construction},
research Notebooks 14--16 and 19--24 in this workspace.
The ordinary reference trajectory, uniform classical minimum,
ordinary representative transfer and strict support calculus
are the principal inputs.
\bibitem{balaban10925}
T.~Ba\l aban,
\emph{Renormalization group approach to lattice gauge field
theories I}, Commun. Math. Phys. \textbf{109} (1987), 249--301,
\href{https://scispace.com/pdf/renormalization-group-approach-to-lattice-gauge-field-3ygai0vxt1.pdf}{original article copy}.
\bibitem{source25}
E.~Sussman and collaborators,
\href{https://github.com/ethanwsussman/Yang--Mills-LLM-experiment}{Yang--Mills-LLM-experiment},
deposited raw accounts at commit
\texttt{e392ef45013a581664bab76b90867170b30fe10e}.
\end{thebibliography}
\end{document}
